\documentclass[11pt,a4paper]{article}
\usepackage[margin=2.5cm]{geometry}
\usepackage{amsmath,amssymb}
\usepackage[T1]{fontenc}
\usepackage{hyperref}
\usepackage{parskip}

\title{\textbf{EB-JEPA for EEG — Concepts and Intuitions}\\
\large A pure-concept summary, no maths}
\author{EB-JEPA EEG track}
\date{June 2026}

\begin{document}
\maketitle

\section*{What is a representation?}

An EEG recording is a sequence of voltages across 19 electrodes over time.
These raw numbers are hard to reason about directly — they vary enormously
across patients, devices, and recording conditions.

A \textbf{representation} is a compact summary of the recording: a vector of,
say, 256 numbers that captures what \emph{matters} about the signal and ignores
what doesn't.  A good representation is one where recordings that are clinically
similar (both abnormal, or both showing the same type of event) end up close to
each other in that 256-dimensional space, and clinically different recordings
end up far apart.

The \textbf{encoder} is the function that maps a raw EEG window to its
representation.  Training an encoder is training it to produce these
meaningful summaries.

\section*{Why self-supervised? The labelling bottleneck}

To train a supervised model you need labelled examples: a neurologist must
inspect each recording and write ``normal'' or ``abnormal''.  This is
expensive, slow, and requires rare expertise.  Most EEG data in the world
is unlabelled.

\textbf{Self-supervised learning (SSL)} sidesteps this: it trains an encoder
to produce good representations \emph{without any labels}, using only the
structure of the raw signal itself.  The idea is to give the encoder a task
that forces it to learn useful structure — not ``is this abnormal?'' but
something like ``do these two corrupted versions of the same signal look like
the same thing?''

Once the encoder is trained on unlabelled data, a very simple classifier
(a single linear layer, called a \emph{probe}) is trained on a small
labelled set.  If the representations are good, the probe can separate
normal from abnormal without ever retraining the encoder.

\section*{The cave allegory}

In Plato's allegory of the cave, prisoners see only shadows projected on a
wall.  They have never seen the real objects casting the shadows, but they
learn to recognise and predict the shadow patterns.

Our SSL encoder is in the same position.  It sees raw EEG voltages — the
shadows.  It never sees the neurologist's verdict — the real object behind
the signal.  But by learning that \emph{two different corrupted views of the
same shadow are shadows of the same object}, it is forced to build an
internal model of what is really there.

A linear probe then asks: ``given your model of what is really there, is
this patient's brain behaving normally?''

\section*{What is invariance? What is collapse?}

\textbf{Invariance} is the property we want: the encoder should output the
same representation for two corrupted versions of the same signal.
Forcing $z_1 \approx z_2$ (where $z_1$ and $z_2$ are the representations
of the two views) is the training signal.

The problem: the trivially simplest way to satisfy $z_1 = z_2$ is to output
the \textbf{same constant vector for everything}.  This is called
\textbf{collapse} — the encoder learns nothing, because all signals look
the same.

\textbf{Anti-collapse regularisation} prevents this.  VICReg does it by
adding two extra penalties: one that forces the encoder to \emph{spread}
its output across dimensions (variance term), and one that forces different
output dimensions to carry different information (covariance term).
SIGReg does it by testing that the distribution of representations looks
Gaussian (a non-constant distribution cannot be collapsed).

\textbf{Can predictive architectures avoid collapse without regularisation?}
There is a common claim that in I-JEPA-style models --- where the encoder
predicts the representation of a \emph{different, disjoint} part of the
input --- collapse is impossible: ``if the encoder outputs a constant, the
MSE to the target cannot be zero.''  This is \emph{not quite right}.  If
both the online and EMA encoders collapse to the same constant $c$, the
predictor can simply learn to output $c$ too, and the loss is zero.

What actually prevents collapse in practice is a combination of:
(1) \textbf{BatchNorm} --- normalisation layers make a true constant output
numerically explosive; (2) \textbf{EMA momentum} --- the target encoder lags
behind the online encoder, creating a non-stationary signal that is hard to
collapse onto; and (3) \textbf{input diversity} from the disjoint split,
which gives useful gradient at the start of training.
The temporal split \emph{helps}, but is not the full story.

\section*{What does the score mean?}

\textbf{BACC} (balanced accuracy) is the average of the true positive rate
(how often abnormal recordings are correctly detected) and the true negative
rate (how often normal recordings are correctly classified).  It is fair even
when the two classes are unequal in size.  Random guessing gives 0.5;
perfect classification gives 1.0.

\textbf{AUROC} (area under the ROC curve) measures how well the model
\emph{ranks} recordings by their abnormality score, at every possible
threshold.  1.0 is perfect; 0.5 is random.

\textbf{Per-window} vs \textbf{per-recording}: a recording is typically
several minutes long, cut into 10-second windows.  Per-window scoring
classifies each window independently; per-recording scoring averages across
all windows of the same patient.  Per-recording scores are systematically
about 5 percentage points higher for the same encoder, because averaging
reduces noise.  Most of the published literature reports per-window;
our primary evaluation reports per-recording (the clinically meaningful metric).
Comparing the two without care manufactures false wins.

\section*{What is transfer?}

TUAB is a binary task (normal / abnormal).
TUEV is a 6-class task (spike-and-slow-wave, generalised periodic discharge,
lateralised periodic discharge, eye movement, artefact, background).

\textbf{Transfer} means: train the encoder on TUAB without labels, then freeze
it and train a linear probe on TUEV.  If the representations are genuinely
general — not just memorised TUAB-specific patterns — they should also separate
TUEV event types, even though the encoder never saw a TUEV label.

A model that transfers well is learning something about the \emph{structure
of EEG} rather than the specific abnormality label.

\section*{Why does the simplest encoder win?}

A 0.4-million-parameter strided convolutional network (Conv1D) consistently
outperforms larger transformer architectures (1.2M--3.65M parameters) in
our experiments.

The intuition: transformers are powerful but data-hungry.  They need millions
of examples and long training to learn useful attention patterns.  On 2\,717
recordings trained for 20 epochs, there is simply not enough signal to fill
their capacity.  The small conv has less capacity but fits the data regime.

The published foundation models (LaBraM, BIOT) that do outperform our conv
were pretrained on tens of thousands of hours of EEG data across multiple
hospitals and devices — a scale we do not have access to.

\end{document}
